\subsection*{The TriLens stack}
The model named in this paper's analyses, TriLens, reads the three genomic views through three
branches and joins them with five classical learners. Its full specification, including every
equation, is in Supplementary Methods; the design rationale and the parts that matter for reading
the results are summarised here.

The motivation is that a resistance mechanism takes more than one genomic form. A point mutation in
a target enzyme is invisible to a presence/absence profile, because the gene is present in
resistant and susceptible strains alike; an acquired plasmid gene is invisible to a mutation-based
view, because it is not a mutation but an arrival. Rather than choose, the model reads all three
views and learns per drug how much each is worth.

\textbf{Three branches.} A \emph{motif encoder} reads the SNP blocks with two stacked
one-dimensional convolutions (width 3, 16 then 32 channels) inside each gene, masked max-pooling
over positions, and gated attention \cite{ilse2018attention} over genes. Convolution is confined to
gene blocks for a measured reason: a plain convolutional network over the whole pan-genome variant
sequence scored 0.775 average precision for ciprofloxacin and 0.632 after the variant order was
shuffled on the same data, so pan-genome position is not meaningless, but within a gene sequence
order is biologically real and that is where the motif encoder works. A \emph{gene attention
module} reads the presence/absence profile, assigns each locus a learned embedding gated by its
presence, and pools with gated attention. A \emph{cross-antibiotic module} reads the per-gene
statistics, pools them, splits the pooled representation into one projection per drug class, and
adds a co-resistance refinement term initialised near zero. Each branch emits a 64-dimensional
summary and seven per-drug scores.

\textbf{Strain conditioning.} On a presence/absence vector, standard gated attention is degenerate:
the scored item for any locus the strain carries is that locus's embedding, identical across
strains, so the score is strain-independent and the softmax merely renormalises over whichever loci
the strain happens to have. The lens then collapses to a global gene ranking and cannot express
that a locus matters for a particular isolate. We condition the score on the mean embedding of the
loci the strain actually carries, which restores per-sample behaviour at the cost of two projection
matrices, and ablate the term in the Results. The attention sum over 32,441 loci is accumulated in
single precision; at half precision the round-off disrupts training.

\textbf{Gate.} The three summaries are concatenated and passed through a small network and a softmax
to give, for each strain and each drug, three weights that sum to one; the fused score is the
weighted sum of the branch scores. A Shannon-entropy term on the gate, with an auxiliary loss on
each branch, is included because without them the gate concentrates on whichever branch is
momentarily ahead and the other two stop receiving gradient.

\textbf{Classical arm and meta-learner.} Five base learners run on the gene features: XGBoost on
gene-aggregated features, XGBoost on a $\chi^2$-prefiltered raw SNP matrix, LightGBM
\cite{ke2017lightgbm}, random forest \cite{breiman2001rf} and logistic regression. They were chosen
to fail differently, which is what a stacking ensemble exploits: two boosting implementations that
grow trees differently, a low-variance averaged forest, a linear model that cannot invent
interactions, and one learner reading single variants rather than gene summaries. Each is trained
in five-fold cross-validation and contributes out-of-fold predictions, so no base prediction is
made by a model that saw the strain. An $L_2$-penalised logistic regression per drug ($C = 1$)
combines the neural probability with the five classical ones \cite{wolpert1992stacked}, and a fixed
threshold of 0.5 turns the result into a call. All seven drugs are predicted in one pass.

A neural meta-learner (one token per drug, four-head self-attention, per-drug decoding heads) was
also tested and scored lower (0.825 against 0.850 average precision on the random split) on about
5,000 training rows, so the linear meta-learner is used and its weights are readable.

Class imbalance is handled by weighting resistant and susceptible samples inversely to frequency
(\texttt{scale\_pos\_weight} in the boosting learners, balanced class weights in the others,
positive class weighting in the neural loss). The neural branches, gate and fusion are trained
jointly with AdamW, learning rate $10^{-3}$ on a cosine schedule, weight decay $10^{-4}$, batch
size 16, gradient clipping at norm 1.0, up to 50 epochs with early stopping after 8 epochs without
improvement in validation mean average precision. Table~\ref{tab:hyper} lists the settings.

\subsection*{The simpler models compared against it}
\textbf{Gradient boosting on determinant calls.} XGBoost \cite{chen2016xgboost} with fixed settings
(400 trees, depth 6, learning rate 0.05, subsample 0.8, column subsample 0.5), no tuning, one model
per drug. ``Acquired-only'' uses the acquired-gene columns; ``all calls'' adds point mutations.

\textbf{Sparse linear model.} $L_1$-penalised logistic regression (scikit-learn
\cite{pedregosa2011sklearn}, liblinear, $C = 1$) on the same binary features, with no tuning.

\textbf{Determinant rule.} A strain is called resistant if it carries a determinant from a fixed
per-drug list (Table~\ref{tab:rule}), written from known mechanisms before any rule score was
computed. It needs no training.

\textbf{Assembly arm.} Histogram gradient boosting on assembly length, contig count and N50 only.

\textbf{External baselines.} TabNet \cite{arik2021tabnet} on the same gene-aggregated features; an
XGBoost model on 1,024 principal components; and a null model using only the eight-way phylogroup
label.

\subsection*{Evaluation splits}
Three splits partition the training collection, at 70/15/15 with a fixed seed (42), and the test
partition is untouched until final evaluation.

The \emph{random split} (4,872/1,045/1,045) stratifies on the ampicillin label. It is the setting
most published work uses, and related strains fall on both sides.

The \emph{phylogroup-stratified split} (same sizes) balances the eight-way phylogroup across
partitions. We include it to make a distinction explicit: every phylogroup appears on both sides,
so it does not test an unseen lineage, only the effect of the stratification variable. We built it
before recognising that it could not answer our question.

The \emph{sequence-type lineage holdout} (4,836/1,081/1,045) assigns whole sequence types. Types
were taken from the metadata's sequence-type field, joined on the sequencing accession carried by
both tables rather than on row position, since joining by row would silently mislabel every
strain's lineage. A type was called for 6,487 strains (93.2\%) across 458 types. Two rules were
fixed before any result was seen: a strain without a type goes to training, as an unlabelled group
cannot be held out safely; and any type holding over 5\% of the collection goes to training, so
that no test set is dominated by one clone. Four types met the second rule (ST73, with 884 strains
or nearly one in eight, ST131, ST95, ST69). The remaining 454 types were allocated whole by a
seeded permutation walk. No sequence type appears in more than one partition, verified rather than
assumed, and no strain without a type is in the test partition. Because whole groups move together,
the partition sizes are not round. The gene selection is refitted on each split's training
partition; the holdout's gene set overlaps the random split's by less than half.

Two further consequences of taking the types as given: the scheme resolves clones but not clonal
complexes, so two single-locus variants can still be separated across partitions and the holdout is
a lower bound on lineage novelty; and any upstream typing error propagates undetected.

For the pre-registered experiments we also used ST-grouped five-fold cross-validation (each type
whole in one fold), random five-fold cross-validation, two country splits (United
Kingdom $\rightarrow$ Norway and the reverse) and a temporal split within Norway (training before
2013, test from 2013).

\subsection*{The external Giessen collection}
A second, independent \Ecoli{} collection from the University of Giessen \cite{ren2022giessen}
provides a cross-collection comparison of the method rather than of a model: Giessen variants are
ordered along a single reference genome while the training collection uses a pan-genome, the two
feature spaces do not overlap, so no model can be carried across. We use the 809 isolates with
complete results for the four shared drugs, described by 60,936 SNP positions, split 647/81/81
stratified on ciprofloxacin. Giessen phenotypes are interpreted against EUCAST breakpoints and the
training phenotypes against a different standard, and breakpoints differ for the same drug, so
label definition and model behaviour cannot be separated in this comparison; prevalence is also two
to three times higher for every shared drug (Table~\ref{tab:cohorts}).

\subsection*{Metrics and uncertainty}
The primary metric for the three internal splits is the area under the precision--recall curve
(average precision), which does not reward majority-class prediction; we also report the Matthews
correlation coefficient, $F_1$, $F_2$ (weighting recall four times precision) and accuracy, each
per drug, at the fixed 0.5 threshold.

Because test prevalence differs between splits, drugs and cohorts, the pre-registered experiments
report prevalence-normalised average precision, $\mathrm{nAP} = (\mathrm{AP} - \pi)/(1 - \pi)$,
where $\pi$ is the test prevalence: 0 is the score of a random ranking and 1 a perfect one.
Within-lineage discrimination is the mean, over sequence types containing both phenotypes, of the
AUROC inside each type; a score that is constant within a type gets exactly 0.5.

Intervals on the three internal splits are 95\% percentile bootstrap intervals over the test set
(1,000 resamples). In the pre-registered experiments they are 95\% paired bootstrap intervals
(2,000 replicates) resampling whole clusters: sequence types internally, BioProjects externally,
SNP clusters in \textit{Campylobacter}. A test-set bootstrap covers which strains fall in the test
partition and not training variability; the seed-stability experiment addresses the latter directly
for the stack's holdout margin, and differences of a few thousandths elsewhere are reported as
directions rather than effects.

\subsection*{Near-duplicate audit}
Bacterial collections contain near-identical strains, and a copy on each side of a split would let
a model score by recognition. We computed the Hamming distance from each test strain to its nearest
training strain. On the random split the closest pair differed at 0.26\% of SNP positions and 47 of
1,045 test strains lay within 0.5\% of a training strain; on the lineage holdout none did. The
0.5\% cut is a working threshold, corresponding to roughly 2,700 differing positions on the
535,439-position matrix, an order of magnitude below the typical separation of two strains of the
same sequence type here. We report the score with those strains removed rather than resting on the
cut. No test genome is identical to a training genome on the filtered SNPs in any of the three
splits.

\subsection*{External cohort}
The NCBI Pathogen Detection \Ecoli{}/\textit{Shigella} isolate table (PDG000000004.6395) was
filtered to isolates with submitted susceptibility results, called with AMRFinderPlus 4.2.7, and
not matching any of our strains by run, BioSample or assembly accession. The primary cohort uses
the same database version as training (10,232 isolates in 62 BioProjects); all 4.2.7 isolates
(10,876) form a sensitivity analysis. Models were trained once on all called training strains and
applied unchanged, with external genes absent from training dropped. CTX was tested on only 1,054
external isolates, 70\% resistant, and is reported separately from the primary mean.

\subsection*{MIC harmonisation}
MICs, measurement signs and testing standards for the external cohort were exported from the
Pathogen Detection AST browser (171,750 rows; all 10,232 isolates matched by PDT accession).
Censored values became intervals ($\le v \mapsto (0, v]$; $< v \mapsto (0, v/2]$;
$\ge v \mapsto [v, \infty)$; $> v \mapsto [2v, \infty)$), and an isolate was labelled with the
EUCAST v16.1 Enterobacterales breakpoints \cite{eucast2026v161} only when the whole interval fell
on one side of a breakpoint. Trimethoprim--sulfamethoxazole values reported as the sulfamethoxazole
component (19 times a trimethoprim dilution) were divided by 19. CTX had fewer than 1,000 MICs and
was not analysed.

\subsection*{Second species}
All 107 organism folders in Pathogen Detection were screened with a pre-registered gate: at least
2,000 isolates with susceptibility results and AMRFinderPlus 4.2.7 calls, at least three drugs with
100 resistant and 100 susceptible isolates, SNP-cluster membership for at least 80\% of isolates
across at least 50 clusters, and assembly statistics for at least 80\%. One organism passed,
\textit{Campylobacter} (4,426 isolates). Its determinant rule (gyrA point mutations for
ciprofloxacin and nalidixic acid; tet genes for tetracycline; 23S rRNA point mutations or erm genes
for clindamycin) was taken from the literature before any score was computed
\cite{zhao2016campylobacter,iovine2013campylobacter}. Folds were grouped by SNP cluster.

\subsection*{Clinical error rates}
Following the FDA guidance for susceptibility-testing systems \cite{fda2009ast}, a major error is a
susceptible isolate called resistant (as a share of susceptible isolates), a very major error a
resistant isolate called susceptible (as a share of resistant isolates), and categorical agreement
the share of correct calls. Learned models were called resistant at a score of 0.5, fixed in
advance. The guidance's benchmarks (categorical agreement $> 89.9\%$, major error $< 3\%$, upper
95\% Clopper--Pearson limit of very major error $< 7.5\%$) are reported as context only: they were
written for phenotypic devices verified against a reference method on a challenge panel enriched
for near-breakpoint isolates, not for retrospective genomic prediction on a convenience collection,
and our two-category comparison is not the quantity the guidance defines. Collapsing intermediate
into resistant can only raise our apparent very major error rate relative to the standard
definition, so the direction of that bias is against us.

To test whether the clinical shortfall is a matter of where the cut sits, we also fitted per-drug
thresholds on the validation partition, once maximising $F_2$ and once constraining the major error
rate, and applied them to the test partition.

\subsection*{Pre-registration}
Each experiment (W6--W14) has a document stating its question, design, endpoint and pass rule,
committed to a public repository with a timestamp before its data were read; later amendments are
dated and were also committed before the data they concern. The three-split evaluation, the
architecture ablations and the Giessen comparison predate this scheme and are reported as
exploratory analyses of a fixed pipeline rather than as pre-registered tests.
Table~\ref{tab:prereg} lists every pre-registered endpoint and its outcome, including those that
failed.
